\documentclass[10pt]{article}
\usepackage[margin=2cm]{geometry}
\usepackage{booktabs}
\usepackage{amsmath}
\usepackage{enumitem}
\usepackage[most]{tcolorbox}
\usepackage{hyperref}
\setlength{\emergencystretch}{3em}
\setlength{\parindent}{0pt}
\setlength{\parskip}{5pt}

% Empty box to write the answer in
\newcommand{\answerbox}[1]{%
  \begin{tcolorbox}[colback=white, colframe=black!40, boxrule=0.4pt,
                    arc=0pt, height=#1]\end{tcolorbox}}

\begin{document}

\begin{center}
{\Large\bfseries CSCI 5942: Mock Quiz 6}\\[3pt]
{\large Assignment 6: Serving GPT-2}
\end{center}

\begin{tcolorbox}[colback=black!4, colframe=black!40, boxrule=0.4pt, arc=0pt]
\begin{itemize}[leftmargin=*, nosep]
\item 15 minutes per set. Closed book, no devices. 3 points per set.
\item Answer as you would in a design interview: two to four sentences per
  part, with the reasoning.
\item Every number you need is on this sheet. Do \textbf{Set A} first and
  check the answers, then do \textbf{Set B}. \textbf{Set C} is about
  the ideas behind your A6 code.
\item 1\,GB $= 10^9$ bytes.
\end{itemize}
\end{tcolorbox}

% =====================================================================
\section*{Set A}

\textbf{The setting.} The restaurant-booking site from Quiz 5 is live.
Every listing page asks the model for one sentence, on demand. You serve
\textbf{GPT-2 large} (774M parameters, \textbf{36 layers, width 1280, 20
heads}, 1{,}024-token context) in \textbf{bf16} with your Assignment 6
engine, on one of the team's 24\,GiB cards. After the weights and the
workspace, \textbf{20\,GB is left for the KV cache}.

A listing request is a \textbf{70-token prompt} (instruction line +
attribute list) and \textbf{at most 40 new tokens}.

\begin{center}
\textbf{Decode step time on this card} (one token for every sequence in the batch)\\[4pt]
\begin{tabular}{lrrrrrr}
\toprule
batch $B$ & 1 & 16 & 64 & 128 & 256 & 512 \\
\midrule
step time (ms) & 14.0 & 14.3 & 15.1 & 17.0 & 24.5 & 41.0 \\
\bottomrule
\end{tabular}
\end{center}

\subsection*{1. Sizing the cache (1.5 points)}

\textbf{(a)} How many bytes of KV cache does one token take? How many
listing requests fit in the 20\,GB at once with a \textbf{contiguous}
cache that reserves the full 1{,}024-token context per request, and how
many with a \textbf{paged} cache at block size 16? Show your working.
\answerbox{4.2cm}

\textbf{(b)} Product sets the TBT (time between tokens) SLO at
\textbf{30\,ms}. What \texttt{max\_batch} would you set, and what
tokens/s do you expect? Compare with the contiguous cache. Product then
says, ``Paging fits about $9\times$ more requests, so we should get about
$9\times$ the throughput.'' What is your response?
\answerbox{4.2cm}

\newpage
\subsection*{2. The engine in production (1.5 points)}

\textbf{(a)} To make room for a second product, the listings get only
\textbf{1\,GB} of KV. You serve \textbf{600 listing requests that arrive
at once} (prompts 60--80 tokens, outputs 10--40 tokens, block size 16,
max batch 128):

\begin{center}
\begin{tabular}{lrrrrrr}
\toprule
engine & time (s) & tok/s & mean lat (s) & p90 lat (s) & mean batch & preempt \\
\midrule
static                  & 7.3 & 2{,}050 & 4.1 & 6.6 & 31.0 & 0 \\
continuous (safe)       & 5.4 & 2{,}780 & 2.9 & 4.9 & 45.2 & 0 \\
continuous (optimistic) & 4.5 & 3{,}350 & 2.4 & 4.2 & 62.5 & 38 \\
\bottomrule
\end{tabular}
\end{center}

Your manager says: ``Optimistic preempted 38 requests. Each of those did
its work twice, so it is wasting GPU time. Turn it off.'' Do you agree?
Use the numbers, and explain why a preemption costs less than it looks.
\answerbox{5.5cm}

\textbf{(b)} A new hire rewrites \texttt{PagedKVCache.update}. Every test
that decodes \textbf{one} sequence passes, and a batch of prompts
\textbf{with the same length} still matches the uncached model. But in a
batch with prompts of \textbf{different lengths}, the longest prompt's
output is correct and the \textbf{shorter prompts produce wrong tokens}.
What is the most likely bug, and why does it only show up for the shorter
prompts? (Hint: think about what the block table holds for a short row.)
\answerbox{5.5cm}

% =====================================================================
\newpage
\section*{Set B}

\textbf{The setting.} Same model and card. The team now also runs the
screenwriting \textbf{autocomplete} on it: many users are decoding at
once, and each one watches the text appear.

\subsection*{3. Sharing and scheduling (1.5 points)}

\textbf{(a)} Every listing prompt is a \textbf{fixed 50-token instruction
line} followed by a 20-token attribute list. You turn on vLLM-style
\textbf{automatic prefix caching} (16-token blocks, each block's hash
includes its parent's hash). How many prompt tokens per request can be
reused, and what fraction of the prompt is that? Then a teammate adds
\texttt{Generated at 2026-11-03 14:02:11} to the \textbf{start} of every
prompt, for logging, and the cache hit rate drops to 0\%. Why, and how
would you fix it?
\answerbox{5cm}

\textbf{(b)} A new feature summarizes a whole screenplay scene:
\textbf{1{,}000-token prompts}. About \textbf{100 autocomplete users} are
decoding at any time. When a summary request arrives, the engine runs its
whole prefill in one step together with the 100 decodes, and every
autocomplete user sees the text stutter. Step time against the number of
tokens in one step:

\begin{center}
\begin{tabular}{lrrrrr}
\toprule
tokens in the step & 64 & 256 & 512 & 768 & 1{,}024 \\
\midrule
step time (ms) & 15 & 18 & 27 & 38 & 50 \\
\bottomrule
\end{tabular}
\end{center}

Which Lecture 15 technique fixes the freeze? With a \textbf{30\,ms TBT
SLO}, what token budget $\tau$ would you pick, and what does the summary
request pay for it?
\answerbox{5cm}

\newpage
\subsection*{4. Design from scratch (1.5 points)}

\textbf{(a)} A customer wants to self-host \textbf{Llama-3.1-8B} (32
layers, \textbf{8 KV heads} $\times$ head dim 128, bf16) for chat with
\textbf{8K-token} conversations, on one \textbf{80\,GB} GPU. The weights
take 16\,GB, and 4\,GB is left for activations. How many 8K conversations
fit at once? Someone proposes switching to an older model with the
\textbf{same size and shape but full multi-head attention (32 KV heads)}.
How many fit then?
\answerbox{6cm}

\textbf{(b)} The customer's users come back to the same conversation many
times a day. Each new turn re-sends the whole history. Name two ways to
avoid recomputing the history's prefill, and say what limits each one.
\answerbox{6cm}

% =====================================================================
\newpage
\section*{Set C: The ideas behind your A6 code}

\textbf{The setting.} No code needed. These questions are about the
design choices inside your A6 cache and scheduler.

\subsection*{5. Inside the cache and the scheduler (1.5 points)}

\textbf{(a)} Your contiguous KV cache runs two prompts in one batch: X has
3 tokens and Y has 5. The cache pads X on the \textbf{left} with 2 empty
slots, so both rows are 5 slots long.
\begin{enumerate}[label=(\roman*), nosep]
\item Why pad on the left and not on the right? (Hint: where does each
  row's next token go during decoding?)
\item What position should X's first real token get, and which slots must
  X's tokens never attend to?
\item A teammate gives every token its slot number as its position, so
  X's tokens get positions 2, 3, 4. Nothing crashes and Y's output is
  correct, but X's output differs from running X alone. Why?
\end{enumerate}
\answerbox{5cm}

\textbf{(b)} Your paged engine uses \textbf{safe admission}: it admits a
request only if the pool can hold the request's \textbf{full length}
(prompt + max new tokens), in blocks of 16. Blocks are handed out as a
sequence grows, not all at once, and a newly admitted request holds no
blocks until its prefill. The pool has \textbf{20 blocks}, and two
requests are running:

\begin{center}
\begin{tabular}{lrr}
\toprule
request & full length (tokens) & blocks held now \\
\midrule
R1 & 100 & 4 \\
R2 & 64 & 2 \\
\bottomrule
\end{tabular}
\end{center}

Two requests wait in this order: A (full length 120 tokens), then B (60
tokens).
\begin{enumerate}[label=(\roman*), nosep]
\item Is A admitted? Then B? Show your working.
\item A teammate simplifies the rule to ``admit if the request's blocks
  fit in the blocks that are free right now.'' Nothing goes wrong at
  admission, but later the engine runs out of blocks. Why later and not
  at admission? Why is running out a bug in safe mode, while optimistic
  mode handles it?
\end{enumerate}
\answerbox{5cm}

\end{document}
